
\documentclass[a4paper, 12pt, leqno]{amsart}
\usepackage{amsmath, mathtools}
\usepackage{amscd}
\usepackage{amssymb}
\usepackage{amsthm}
\usepackage{ascmac, color}
\usepackage{graphicx, xcolor, pdfpages}
\usepackage{mathrsfs}
%\usepackage[all]{xy}
%\usepackage[notcite,notref]{showkeys}
\usepackage{hyperref}


\setlength{\textwidth}{15.0truecm}
\setlength{\textheight}{22.5truecm}
\setlength{\topmargin}{0mm}
\setlength{\oddsidemargin}{0.3cm} 
\setlength{\evensidemargin}{0.3cm}



\newtheorem{lemma}{{\sc Lemma}}[section]
\newtheorem{corollary}[lemma]{{\sc Corollary}}
\newtheorem{proposition}[lemma]{{\sc Proposition}}
\newtheorem{theorem}[lemma]{{\sc Theorem}}
\theoremstyle{definition}
\newtheorem{remark}[lemma]{{\sc Remark}}
\newtheorem{definition}[lemma]{{\sc Definition}}
\newtheorem{example}[lemma]{{\sc Example}}
\newtheorem{conjecture}[lemma]{{\sc Conjecture}}
\newtheorem{assumption}[lemma]{{\sc Assumption}}
\renewcommand{\proofname}{{\sc Proof}}
\numberwithin{equation}{section}
\renewcommand{\theequation}{\thesection.\arabic{equation}}
\renewcommand{\thelemma}{\thesection.\arabic{lemma}}



\def\Ga{{\mathfrak{a}}}
\def\Gb{{\mathfrak{b}}}
\def\Gc{{\mathfrak{c}}}
\def\Bd{{\mathfrak{d}}}
\def\Gg{{\mathfrak{g}}}
\def\Gh{{\mathfrak{h}}}
\def\Gk{{\mathfrak{k}}}
\def\Gl{{\mathfrak{l}}}
\def\Gm{{\mathfrak{m}}}
\def\Gn{{\mathfrak{n}}}
\def\Go{{\mathfrak{o}}}
\def\Gp{{\mathfrak{p}}}
\def\Gr{{\mathfrak{r}}}
\def\Gs{{\mathfrak{s}}}
\def\Gt{{\mathfrak{t}}}
\def\Gz{{\mathfrak{z}}}
\def\Gsl{{\mathfrak{sl}}}
\def\Gsp{{\mathfrak{sp}}}
\def\GA{{\mathfrak{A}}}
\def\GC{{\mathfrak{C}}}
\def\GD{{\mathfrak{D}}}
\def\GE{{\mathfrak{E}}}
\def\GF{{\mathfrak{F}}}
\def\GH{{\mathfrak{H}}}
\def\GI{{\mathfrak{I}}}
\def\GJ{{\mathfrak{J}}}
\def\GK{{\mathfrak{K}}}
\def\GL{{\mathfrak{L}}}
\def\GM{{\mathfrak{M}}}
\def\GN{{\mathfrak{N}}}
\def\GO{{\mathfrak{O}}}
\def\GS{{\mathfrak{S}}}
\def\GU{{\mathfrak{U}}}
\def\GV{{\mathfrak{V}}}
\def\GX{{\mathfrak{X}}}
\def\GZ{{\mathfrak{Z}}}



\def\BA{{\mathbf{A}}}
\def\BB{{\mathbf{B}}}
\def\BC{{\mathbf{C}}}
\def\BD{{\mathbf{D}}}
\def\BE{{\mathbf{E}}}
\def\BF{{\mathbf{F}}}
\def\BG{{\mathbf{G}}}
\def\BI{{\mathbf{I}}}
\def\BJ{{\mathbf{J}}}
\def\BK{{\mathbf{K}}}
\def\BN{{\mathbf{N}}}
\def\BP{{\mathbf{P}}}
\def\BQ{{\mathbf{Q}}}
\def\BR{{\mathbf{R}}}
\def\BX{{\mathbf{X}}}
\def\BZ{{\mathbf{Z}}}

\def\Ba{{\mathbf{a}}}
\def\Bd{{\mathbf{d}}}
\def\Bl{{\mathbf{l}}}
\def\Bm{{\mathbf{m}}}
\def\Bs{{\mathbf{s}}}
\def\Bt{{\mathbf{t}}}

\def\CA{{\mathcal A}}
\def\CB{{\mathcal B}}
\def\CC{{\mathcal C}}
\def\DD{{\mathcal D}}
\def\CE{{\mathcal E}}
\def\CF{{\mathcal F}}
\def\CH{{\mathcal H}}
\def\CG{{\mathcal G}}
\def\CO{{\mathcal O}}
\def\CI{{\mathcal I}}
\def\CJ{{\mathcal J}}
\def\CK{{\mathcal K}}
\def\CL{{\mathcal L}}
\def\CM{{\mathcal M}}
\def\CN{{\mathcal N}}
\def\CP{{\mathcal P}}
\def\CQ{{\mathcal Q}}
\def\CR{{\mathcal R}}
\def\CS{{\mathcal S}}
\def\CT{{\mathcal T}}
\def\CU{{\mathcal U}}
\def\CV{{\mathcal V}}
\def\CW{{\mathcal W}}
\def\CX{{\mathcal X}}
\def\CY{{\mathcal Y}}
\def\CZ{{\mathcal Z}}

\def\SI{{\mathscr I}}


\def\Ad{\mathop{\rm Ad}\nolimits}
\def\ad{{\mathop{\rm ad}\nolimits}}
\def\alg{{\mathop{\rm alg}\nolimits}}
\def\an{{\mathop{\rm an}\nolimits}}
\def\co{\mathop{\rm co}\nolimits}
\def\Cok{\mathop{\rm Cok}\nolimits}
\def\Comod{\mathop{\rm Comod}\nolimits}
\def\deru{\partial}
\def\diag{{\mathop{\rm diag}\nolimits}}
\def\End{\mathop{\rm{End}}\nolimits}
\def\EEnd{\mathcal{E}nd}
\def\eq{\mathop{\rm eq}\nolimits}
\def\Ext{{\mathop{\rm Ext}\nolimits}}
\def\fin{\rm fin}
\def\For{{\mathop{\rm For}\nolimits}}
\def\gr{\mathop{\rm gr}\nolimits}
\def\Hom{\mathop{\rm Hom}\nolimits}
\def\hotimes{{\;\mathop{\widehat{\otimes}}\nolimits\;}}
\def\HHom{\mathcal{H}om}
\def\Ht{\mathop{\rm ht}\nolimits}
\def\id{\mathop{\rm id}\nolimits}
\def\Id{\mathop{\rm Id}\nolimits}
\def\Ind{{\mathop{\rm Ind}\nolimits}}
\def\inte{{\mathop{\rm int}\nolimits}}
\def\Image{\mathop{\rm Im}\nolimits}
\def\Ker{\mathop{\rm Ker\hskip.5pt}\nolimits}
\def\Mod{\mathop{\rm Mod}\nolimits}
\def\nil{{\rm{nil}}}
\def\op{{\mathop{\rm op}\nolimits}}
\def\ord{{\mathop{\rm ord}\nolimits}}
\def\Proj{\mathop{\rm Proj}\nolimits}
\def\rad{{\mathop{\rm rad}\nolimits}}
\def\rank{{\mathop{\rm rank}\nolimits}}
\def\Rep{{\mathop{\rm Rep}\nolimits}}
\def\Res{{\mathop{\rm Res}\nolimits}}
\def\res{{\mathop{\rm res}\nolimits}}
\def\sgn{\mathop{\rm sgn}\nolimits}
\def\Spec{{\rm{Spec}}}
\def\Supp{{\rm{Supp}}}
\def\Tor{{\rm{Tor}}}
\def\wt{\mathop{\rm wt}\nolimits}
\def\CEnd{\mathop{{\mathcal{E}}nd}\nolimits}

\def\Dqab{D_q^{ab}}
\def\BDqab{{\BD_q^{ab}}}
\def\zDq{{}^0\!D_q}

\def\m{\Gm_{J,q}}
\def\K{U_q(\Gl_J)}
\def\F{O_q(G/L_J)}
\def\al{\blacktriangleleft}
\def\ar{\blacktriangleright}



%\setcounter{tocdepth}{1}
\begin{document}
\title[Differential calculus on quantized irreducible flag manifolds]
{
Differential calculus on quantized irreducible flag manifolds}
\author{Toshiyuki TANISAKI}
%\address{}
%\email{}
\subjclass[2020]{Primary: 17B37, Secondary: 81R60}
\begin{abstract}
We give a new description of the differential calculus on the quantized irreducible flag manifold  given by Heckenberger-Kolb.
We also define a quantum analogue of the Dolbeault complex associated to an equivariant vector bundle on the irreducible flag manifold, and show that its cohomology groups coincide with the cohomology groups given by Andersen-Polo-Wen defined using the induction functor.
\end{abstract}
\maketitle
%\tableofcontents

\section{Introduction}
\subsection{}
For a smooth complex algebraic variety $Y$ let $Y_\an$ be the associated complex manifold.
For an algebraic  vector bundle $\CL$ on $Y$ let $\CL_\an$ be the corresponding holomorphic vector bundle on $Y_\an$.
Assume that $Y$ is projective.
By Serre's GAGA \cite{Se} the cohomology group $H^r(Y;\CL)$ is naturally isomorphic to the Dolbeault cohomology group $H^{0,r}(Y_\an;\CL_\an)$.
It is an interesting theme to pursuit non-commutative analogues of this fact  (see 
\cite{BeS}).
In this paper we are concerned with the special case of this problem related to quantum groups.

Let $G$ be a connected, simply-connected simple algebraic group over the complex number field $\BC$.
Let $P$ be a parabolic subgroup of $G$, and let $L$ be a Levi subgroup of $P$.
The homogeneous space $X=G/P$  is a smooth projective variety called a (generalized) flag manifold for $G$.
For a finite-dimensional $P$-module $V$ there is associated an algebraic $G$-equivariant vector bundle $\CL_V$ on $X$, and 
we have an isomorphism 
\begin{equation}
\label{eq:isom}
H^r(X;\CL_V)\cong H^{0,r}(X_\an;\CL_{V,\an})
\end{equation}
of $G$-modules.

\subsection{}
We consider the quantum analogue of \eqref{eq:isom} using the quantum groups.
In this paper we assume that  the parameter $q$ for the quantum group is not a root of 1.
Let $O_q(G)$, $O_q(P)$ and $O_q(L)$ be the quantized coordinate algebras of $G$, $P$ and $L$ respectively.
They are Hopf algebras over $\BC$.
A quantum analogue of $H^r(V;\CL_V)$ is naturally defined using the 
derived induction functor 
\[
R^r\Ind:\{\text{$O_q(P)$-comodules}\}\to
\{\text{$O_q(G)$-comodules}\}
\]
(see \cite{APW}).
In fact $R^r\Ind(V_q)$ for an $O_q(P)$-comodule $V_q$ plays the role of $H^r(V;\CL_V)$ as above.
However, it is not so clear what is the counterpart of $H^{0,r}(X_\an;\CL_{V,\an})$ in the quantum setting.
Take a maximal compact subgroup $G_u$ of $G_\an$ such that $L_u=G_u\cap P_\an$ is a maximal compact subgroup of  $L_\an$.
As a real manifold we have $X_\an=G_u/L_u$.
Hence the subalgebra $O_q(G/L)$ of $O_q(G)$ consisting of $O_q(L)$-covariant vectors is regarded as a quantum analogue of the $G_u$-finite part of the algebra of 
$\BC$-valued $C^\infty$-functions on $X_{\an}$.





Assume that the unipotent radical of $P$ is commutative 
(in this case $G/P$ is called an irreducible flag manifold in some referrences).
In \cite{HK1}, \cite{HK2}, \cite{HK3} Heceknberger-Kolb  introduced a quantum analogue of (the $G_u$-finite part of) the Dolbeault complex
\begin{equation}
\label{eq:Dol1}
\Omega^{0,\bullet}_{X_\an}=
[0\to \Omega^{0,0}_{X_\an}\xrightarrow{\overline{\deru}}\Omega_{X_\an}^{0,1}\xrightarrow{\overline{\deru}}\cdots].
\end{equation}
We can regard $O_q(G/L)$ equipped with this analogue of \eqref{eq:Dol1} a quantum analogue of the complex manifold $X_\an$.
Moreover, Carotenutto-D\'{i}az Garc\'{i}a-\'{O} Buachalla \cite{CDO} defined, for a one-dimensional $O_q(L)$-comodule $V_q$, a quantum analogue of (the $G_u$-finite part of) the Dolbeault complex
\begin{equation}
\label{eq:Dol2}
\Omega^{0,\bullet}_{X_\an} \otimes_{O_{X_\an}} \CL_{V,\an}=
[0\to 
\Omega^{0,0}_{X_\an}\otimes_{O_{X_\an}}  \CL_{V,\an}
\xrightarrow{\nabla}
\Omega_{X_\an}^{0,1}\otimes_{O_{X_\an}}  \CL_{V,\an}
\xrightarrow{\nabla}\cdots]
\end{equation}
associated to a line bundle $\CL_{V,\an}$, and proved that the Borel-Weil type theorem holds for the quantum analogue of $H^{0,0}(X_\an,\CL_{V,\an})$ (in the case $V_q$ is one-dimensional).

\subsection{}
In the present paper still assuming that the unipotent radical of $P$ is commutative, we 
give a new description of the quantum analogue of \eqref{eq:Dol1}.
We also show that a quantum analogue of \eqref{eq:Dol2} is defined 
for any $O_q(P)$-comodule not necessarily one-dimensional.
Finally, we show that the quantum analogue of \eqref{eq:isom} holds for any $O_q(P)$-comodule $V_q$.
In particular,
the Borel-Weil-Bott type theorem holds for the quantum analogue of $H^{0,r}(X_\an,\CL_{V,\an})$ for any irreducible $O_q(P)$-comodule $V_q$.

\subsection{}
For a Hopf algebra $H$ over $\BC$ we denote the comultiplication, the counit and the antipode by
\[
\Delta:H\to H\otimes H, 
\qquad
\varepsilon:H\to\BC,
\qquad
S:H\to H
\]
respectively.
For $h\in H$ we set $h^+=h-\varepsilon(h)1$.
For a subalgebra $K$ of $H$ we set
\[
K^+=\{k\in K\mid \varepsilon(k)=0\}.
\]
We use Sweedler's notation 
\[
\Delta^{(n)}(h)=\sum h_{(0)}\otimes\cdots\otimes h_{(n)}
\]
for the iterated comultiplication $\Delta^{(n)}:H\to H^{\otimes n+1}$.



\section{Quantum groups}
\subsection{}
Let $G$ be a connected, simply-connected simple algebraic group over the complex number field $\BC$, and let $H$ be a maximal torus of $G$.
We denote by $\Gg$ and $\Gh$ the Lie algebras of $G$ and $H$ respectively.
Let $\Delta\subset\Gh^*$ be the set of roots for $(\Gg,\Gh)$.
We fix a set of simple roots $\{\alpha_i\mid i\in I\}$, and denote by $\Delta^+$ the corresponding set of positive roots.
We denote by $W\subset GL(\Gh^*)$ the Weyl group.
It is a Coxeter group generated by simple reflections $s_i$ ($i\in I$).
We denote by $w_0$ the longest element of $W$.
Let 
\begin{equation}
\label{eq:bil}
(\;,\;):\Gh^*\times\Gh^*\to\BC
\end{equation}
be the $W$-invariant symmetric bilinear form such that $(\alpha,\alpha)=2$ for short roots.
For $\alpha\in\Delta$ the corresponding coroot $\alpha^\vee$ is defined by $\alpha^\vee=2\alpha/(\alpha,\alpha)^{-1}$.
Let $Q$ 
%and $Q^\vee$ 
be the root lattice 
%and the coroot lattice respectively, 
and let $\Lambda\subset\Gh^*$ be the set of integral weights.
We have $Q\subset \Lambda$.
Moreover, 
by our choice of \eqref{eq:bil}
we have 
%$Q\subset Q^\vee$ and 
$(\Lambda, Q)\subset\BZ$.
We set
\[
Q^+=\sum_{i\in I}\BZ_{\geqq0}\alpha_i.
%\qquad
%\Lambda^+=
%\{\lambda\in\Lambda\mid
%(\lambda,\alpha^\vee)\geqq0\;\text{ for } \alpha\in\Delta^+\}.
\]
For $\alpha\in\Delta$ we denote by $\Gg_\alpha$ the root subspace.
We define sublagebras $\Gn$, $\Gb$ of $\Gg$ by
\[
\Gn=\bigoplus_{\alpha\in\Delta^+}\Gg_\alpha,
\qquad
\Gb=\Gh\oplus\Gn.
\]


Let  $J$ be a subset of $I$. 
We set
\[
\Delta_J=\Delta\cap\left(
\sum_{j\in J}\BZ\alpha_j\right),
\qquad
\Delta_J^+=\Delta^+\cap\Delta_J,
\qquad
\Xi_J=\Delta^+\setminus\Delta_J^+.
\]
We denote the longest element of $W_J=\langle s_j\mid j\in J\rangle$ by $w_J$.
We define subalgebras $\Gl_J$, $\Gb_J$, $\Gm_J$, $\Gp_J$ of $\Gg$ by
\[
\Gl_J=\Gh\oplus\left(\bigoplus_{\alpha\in\Delta_J}\Gg_\alpha\right),
\quad
\Gb_J=\Gh\oplus\left(\bigoplus_{\alpha\in\Delta^+_J}\Gg_\alpha\right),
\quad
\Gm_J=\bigoplus_{\alpha\in\Xi_J}\Gg_\alpha,
\quad
\Gp_J=\Gl_J\oplus\Gm_J.
\]
We denote by $L_J$, $P_J$ the closed connected subgroups of $G$ with Lie algebras $\Gl_J$, $\Gp_J$ respectively.
Let $\sum_{i\in I}m_i\alpha_i$ be the highest root.
The following fact is well-known.
\begin{lemma}
\label{lem:irred}
The following conditions for a subset $J$ of $I$ are equivalent to each other.
\begin{itemize}
\item[(a)] 
$\Gm_J\ne\{0\}$ and $[\Gm_J,\Gm_J]=\{0\}$.
\item[(b)] 
$\Gm_J$ is an irreducible $\Gl_J$-module with respect to the adjoint action.
\item[(c)] 
$J=I\setminus\{i_0\}$ with $m_{i_0}=1$.

\end{itemize}
\end{lemma}



\subsection{}
We fix $q\in\BC^\times$.
In this paper we assume that $q$ is not a root of $1$.
We denote by $U_q(\Gg)$ the quantized enveloping algebra of $\Gg$.
It is a $\BC$-algebra generated by the elements 
\[
e_i, \; f_i\quad(i\in I),\qquad k_\gamma\quad(\gamma\in Q)
\]
satisfying the standard relations.
Its Hopf algebra structure is given by 
\[
\Delta(e_i)=e_i\otimes k_i+1\otimes e_i,
\quad
\Delta(f_i)=f_i\otimes 1+k_i^{-1}\otimes f_i,
\quad
\Delta(k_\gamma)=k_\gamma\otimes k_\gamma,
\]
\[
\varepsilon(e_i)=\varepsilon(f_i)=0,\quad\varepsilon(k_\gamma)=1,
\]
\[
S(e_i)=-e_ik_i^{-1},\quad
S(f_i)=-k_if_i,\quad
S(k_\gamma)=k_{-\gamma},
\]
where $k_i=k_{\alpha_i}$.
We define subalgebras $U_q(\Gh)$, $U_q(\Gn)$, $U_q(\Gb)$ of $U_q(\Gg)$ by
\[
U_q(\Gh)=\langle k_\gamma\mid \gamma\in Q\rangle,
\qquad
U_q(\Gn)=\langle e_i\mid i\in I\rangle,
\qquad
U_q(\Gb)=\langle U_q(\Gh), U_q(\Gn)\rangle.
\]
Then $U_q(\Gh)$, $U_q(\Gb)$ are  Hopf subalgebras of $U_q(\Gg)$.
For $\lambda\in\Lambda$ we define a character $\chi_\lambda:U_q(\Gh)\to \BC$ by 
\[
\chi_\lambda(k_\gamma)=q^{(\lambda,\gamma)}
\qquad(\gamma\in Q).
\]
For a left (resp.\ right) $U_q(\Gh)$-module $M$ and $\lambda\in\Lambda$ we set 
\[
M_\lambda=\{m\in M\mid hm=\chi_\lambda(h)m\;
(\text{resp.} \;
mh=\chi_\lambda(h)m)\;\;
\text{for } h\in U_q(\Gh)\}.
\]
We say that a left or right $U_q(\Gh)$-module is a weight module if we have the weight space decomposition
\[
M=\bigoplus_{\lambda\in\Lambda}M_\lambda.
\]


We define the adjoint action of $U_q(\Gg)$ on $U_q(\Gg)$ by
\[
\ad(u)(x)=\sum u_{(1)}x(S^{-1}u_{(0)})
\qquad (u, x\in U_q(\Gg)).
\]
The subalgebra $U_q(\Gn)$ is a left $U_q(\Gh)$-module with respect to the adjoint action with the weight space decomposition 
\[
U_q(\Gn)=\bigoplus_{\gamma\in Q^+}U_q(\Gn)_\gamma.
\]





Following Lusztig \cite{Lb} we define an algebra automorphism $T_i$ ($i\in I$) of $U_q(\Gg)$ by
\begin{align*}
T_i(e_j)=&
\begin{cases}
-f_ik_i\quad &(i=j)
\\
\sum_{r=0}^{-a_{ij}}(-q_i)^{-r}e_i^{(-a_{ij}-r)}e_je_i^{(r)}
\quad &(i\ne j),
\end{cases}
\\
T_i(f_j)=&
\begin{cases}
-k_i^{-1}e_i\quad &(i=j)
\\
\sum_{r=0}^{-a_{ij}}(-q_i)^{r}f_i^{(r)}f_jf_i^{(-a_{ij}-r)}
\quad &(i\ne j),
\end{cases}
\\
T_i(k_\lambda)=&k_{s_i\lambda}.
\end{align*}
Here, $q_i=q^{(\alpha_i,\alpha_i)/2}$, and 
\[
e_i^{(r)}=
\left(\prod_{s=1}^r\frac{q_i^s-q_i^{-s}}{q-q^{-1}}\right)^{-1}
e_i^r,
\qquad
f_i^{(r)}=
\left(\prod_{s=1}^r\frac{q_i^s-q_i^{-s}}{q-q^{-1}}\right)^{-1}
f_i^r.
\]
Its inverse is given by
\begin{align*}
T_i^{-1}(e_j)=&
\begin{cases}
-k_i^{-1}f_i\quad &(i=j)
\\
\sum_{r=0}^{-a_{ij}}(-q_i)^{-r}e_i^{(r)}e_je_i^{(-a_{ij}-r)}
\quad &(i\ne j),
\end{cases}
\\
T_i^{-1}(f_j)=&
\begin{cases}
-k_i^{-1}e_i\quad &(i=j)
\\
\sum_{r=0}^{-a_{ij}}(-q_i)^{r}f_i^{(-a_{ij}-r)}f_jf_i^{(r)}
\quad &(i\ne j),
\end{cases}
\\
T_i^{-1}(k_\lambda)=&k_{s_i\lambda}.
\end{align*}
For $w\in W$ choose a reduced expression 
$w=s_{i_1}\dots s_{i_r}$
of $w$ and define an algebra automorophism $T_w$ of $U_q(\Gg)$ by 
$
T_w=T_{i_1}\dots T_{i_r}.
$
It does not depend on the choice of a reduced expression.

Let $J\subset I$.
We define subalgebras $U_q(\Gl_J)$, $U_q(\Gb_J)$, $U_q(\Gp_J)$, $U_q(\Gm_J)$ of $U_q(\Gg)$ by
\[
U_q(\Gl_J)=
\langle
U_q(\Gh), e_j, f_j\mid j\in J\rangle,
\qquad
U_q(\Gb_J)=
\langle
U_q(\Gh), e_j\mid j\in J\rangle,
\]
\[
U_q(\Gp_J)=
\langle
U_q(\Gh), e_i, f_j\mid i\in I, j\in J\rangle,
\qquad
U_q(\Gm_J)=U_q(\Gn)\cap T_{w_J}(U_q(\Gn)).
\]
Then $U_q(\Gl_J)$ and $U_q(\Gp_J)$ are Hopf subalgebras of $U_q(\Gg)$.
We have the following (see for example \cite[Lemma 2.8]{TM}).
\begin{lemma}
\label{lem:DUm}
$\Delta(U_q(\Gm_J))\subset U_q(\Gm_J)\otimes U_q(\Gb)$.
\end{lemma}


Similarly to \cite{KMT} 
we have the following.
\begin{lemma}
\label{lem:adUm}
$\ad(U_q(\Gl_J))(U_q(\Gm_J))\subset U_q(\Gm_J)$.
\end{lemma}




We choose a reduced expression 
$w_Jw_0=s_{i_1}\dots s_{i_d}$ of $w_Jw_0$ and set
\begin{equation}
\label{eq:beta}
\beta_r=s_{i_1}\dots s_{i_{r-1}}(\alpha_{i_r})
\qquad(r=1,\dots d).
\end{equation}
Then we have $\Xi_J=\{\beta_1,\dots, \beta_d\}$.
Set 
\begin{equation}
\label{eq:b}
e_{\beta_r}=T_{i_1}^{-1}\dots T_{i_{r-1}}^{-1}(e_{i_r})
\qquad(r=1,\dots d).
\end{equation}
Then we have the following (see \cite[Proposition 40.2.1]{Lb}).
\begin{lemma}
\label{lem:basis}
The set
\[
\{
e_{\beta_1}^{n_1}\dots e_{\beta_d}^{n_d}\mid
(n_1,\dots,n_d)\in(\BZ_{\geqq0})^d\}
\]
forms a basis of $U_q(\Gm_J)$.
\end{lemma}
\subsection{}
For $\Ga=\Gg$, $\Gh$, $\Gp_J$, $\Gl_J$ we say that a left (resp.\ right) $U_q(\Ga)$-module $M$ is integrable if it is a weight module as a $U_q(\Gh)$-module, 
and if we have $\dim U_q(\Ga)m<\infty$
(resp.\ $\dim mU_q(\Ga)<\infty$)
for any $m\in M$ .
%We denote by $\Mod_\inte(U_q(\Ga))$ 
%(resp.\ $\Mod^r_\inte(U_q(\Ga))$)
%the category of integrable left (resp.\ right)
%$U_q(\Ga)$-modules.
Let $A$ be the closed connected subgroup of $G$ with Lie algebra $\Ga$.
We denote by $O_q(A)$ the subspace of $U_q(\Ga)^*$ spanned by the matrix coefficients of  integrable left $U_q(\Ga)$-modules 
(it coincides with the subspace of spanned by the matrix coefficients of  integrable right $U_q(\Ga)$-modules).
Then $O_q(A)$ turns out to be a Hopf algebra whose multiplication, comultiplication, unit, counit, antipode are given by the transpose of the 
comultiplication, multiplication, counit, unit, antipode of $U_q(\Ga)$ respectively.
The canonical pairing
\[
\langle\;,\;\rangle:O_q(A)\times U_q(\Ga)\to\BC
\]
is non-degenerate in the sense that 
$u\in U_q(\Ga)$ satisfying $\langle O_q(A),u\rangle=\{0\}$ is zero, and
$\varphi\in O_q(A)$ satisfying $\langle \varphi, U_q(\Ga)\rangle=\{0\}$ is zero.
We have a $U_q(\Ga)$-bimodule structure of $O_q(A)$ given by
\[
\langle u_1\varphi u_2,u\rangle=
\langle\varphi,u_2uu_1\rangle
\qquad(\varphi\in O_q(A), u_1, u_2, u\in U_q(\Ga)).
\]
The category of integrable left (resp.\ right) $U_q(\Ga)$-modules is naturally equivalent to the category of 
right (resp.\ left) $O_q(A)$-comodules $\Mod^{O_q(A)}$ (resp.\ ${}^{O_q(A)}\!\Mod$). 


\section{Differential calculus}
\subsection{}
In the rest of this paper we fix a subset $J$ of $I$ satisfying the equivalent conditions in Lemma \ref{lem:irred}.
In particular, we have $J=I\setminus\{i_0\}$ for some $i_0\in I$.
Define a numbering $\Xi_j=\{\beta_1,\dots, \beta_d\}$ of elements of $\Xi_J$ as in \eqref{eq:beta}.
For each $\beta\in\Xi_J$ we have $e_\beta\in U_q(\Gm_J)$ given  by \eqref{eq:b}.
We set 
\begin{equation}
\m=\sum_{\beta\in\Xi_J}\BC e_\beta
\subset U_q(\Gm_J).
\end{equation}


Consider the grading 
$
U_q(\Gm_J)=\bigoplus_{n=0}^\infty U_q(\Gm_J)(n)
$
of the algebra $U_q(\Gm_J)$, where $U_q(\Gm_J)(n)$
is the sum of the weight spaces whose weight belongs  to 
$n\alpha_{i_0}+\sum_{j\in J}\BZ_{\geqq0}\alpha_j$.
By the condition (c) in Lemma \ref{lem:irred} and Lemma \ref{lem:basis}
we have 
$U_q(\Gm_J)(0)=\BC$,
$U_q(\Gm_J)(1)=\m$, and $U_q(\Gm_J)(n)$ coincides with the image of 
$\m^{\otimes n}\to U_q(\Gm_J)$ defined by the multiplication of $U_q(\Gm_J)$.
Considering the weights we obtain from Lemma \ref{lem:DUm}, Lemma \ref{lem:adUm} the following.


\begin{lemma}
\label{lem:Dm}
For $x\in\m$ we have
\[
\Delta(x)\in 1\otimes x+\m\otimes U_q(\Gb_J).
\]
\end{lemma}



\begin{lemma}
\label{lem:adm}
$\ad(\K)(\m)\subset\m$.
\end{lemma}




By Lemma \ref{lem:basis} $U_q(\Gm_J)(2)$ has a basis,
\[
\{e_{\beta_m}e_{\beta_n}\mid 1\leqq m\leqq n\leqq d\}.
\]
Hence for $d\geqq r>s\geqq 1$ we can write 
\begin{equation}
\label{eq:rel}
e_{\beta_r}e_{\beta_s}=
\sum_{1\leqq m\leqq n\leqq d}C^{r,s}_{m,n}e_{\beta_m}e_{\beta_n}
\end{equation}
for some $C^{r,s}_{m,n}\in\BC$.
By Lemma \ref{lem:basis} we see that the algebra $U_q(\Gm_J)$ is generated by the elements $e_\beta$ ($\beta\in\Xi_J$) satisfying the defining equations \eqref{eq:rel} for $d\geqq r>s\geqq 1$.
In particular, $U_q(\Gm_J)$ is a quadratic algebra generated by $\m$.
We set
\begin{align}
\Gr=&
\{\sum_c x_c\otimes y_c\in\m\otimes\m\mid
\sum x_cy_c=0\}
\\
\nonumber
=&\sum_{r>s}\BC(e_{\beta_r}\otimes e_{\beta_s}-
\sum_{m\leqq n}C^{r,s}_{m,n}e_{\beta_m}\otimes e_{\beta_n})
\subset \m\otimes\m.
\end{align}

We denote by $\wedge\m^*$ the quadratic algebra dual to $U_q(\Gm_J)$.
Namely, we set
\[
\wedge\m^*=T\m^*/(T\m^*)\Gr^\perp (T\m^*)
\]
where $T\m^*$ is the tensor algebra of the $\BC$-module $\m^*$ and 
\[
\Gr^\perp=\{
z\in\m^*\otimes\m^*\mid \langle z,\Gr\rangle=\{0\}\}.
\]
The multiplication of $\wedge\m^*$ is written as 
\[
\wedge\m^*\times \wedge\m^*\to \wedge\m^*
\qquad((a,b)\mapsto a\wedge b).
\]
For $n\geqq0$ we denote by $\wedge^n\m^*$ the image of $(\m^*)^{\otimes n}\to \wedge\m^*$ defined by the multiplication of $\wedge\m^*$.
Then we have 
\[
\wedge\m^*=\bigoplus_n\wedge^n\m^*.
\]
Denote by $\{e^*_{\beta_1},\dots, e^*_{\beta_d}\}$ the basis of $\m^*$ dual to $\{e_{\beta_1},\dots, e_{\beta_d}\}$.
By a general fact on quadratic algebras (see \cite[Theorem 4.1]{PP}) we have 
$\wedge^n\m^*=\{0\}$ for $n>d(d-1)/2$, and 
$\wedge^n\m^*$ for $0\leqq n\leqq d(d-1)/2$ has a basis
\[
\{e^*_{\beta_{r_1}}\wedge\dots\wedge e^*_{\beta_{r_n}}\mid
d\geqq r_1>\dots>r_n\geqq 1\}
\]
\subsection{}
Since the pairing $O_q(G)\times U_q(\Gg)\to\BC$ is non-degenerate and since $\m$ is a finite-dimensional subspace of $U_q(\Gg)$, the natural linear map $O_q(G)\to\m^*$ given by the restriction is surjective.
We see by Lemma \ref{lem:Dm} that the kernel $\SI$ of $O_q(G)\to\m^*$ is a right ideal of $O_q(G)$.
We regard $\m^*$ as a right $O_q(G)$-module by identifying it with $O_q(G)/\SI$.
We can show the following using Lemma \ref{lem:Dm}.


\begin{lemma}
\label{lem:mO}
The right $O_q(G)$-module structure of $\m^*$ is naturally extended to that of $\wedge\m^*$ by
\begin{itemize}
\item[(a)]
$1\varphi=\varepsilon(\varphi)1$\quad$(\varphi\in O_q(G))$,
\item[(b)]
$(a\wedge b)\varphi=\sum a\varphi_{(0)}\wedge b\varphi_{(1)}$
\quad
$(a, b\in\wedge\m^*, \varphi\in O_q(G))$.

\end{itemize}
\end{lemma}
Hence we obtain a graded algebra structure of $O_q(G)\otimes \wedge\m^*=\bigoplus_{n=0}^{d(d-1)/2} O_q(G)\otimes \wedge^n\m^*$ given by
\begin{equation}
(\varphi\otimes a)\wedge(\psi\otimes b)
=
\sum \varphi\psi_{(0)}\otimes (a\psi_{(1)})\wedge b
\qquad(\varphi, \psi\in O_q(G), a, b\in\wedge\m^*).
\end{equation}
Note that $O_q(G)$ and $\wedge\m^*$ are naturally identified with sublagebras of $O_q(G)\otimes \wedge\m^*$ by the embeddings
$O_q(G)\cong O_q(G)\otimes 1\subset O_q(G)\otimes \wedge\m^*$ and 
$\wedge\m^*\cong1\otimes \wedge\m^*\subset O_q(G)\otimes \wedge\m^*$.
In particular, $O_q(G)\otimes \wedge\m^*$ is naturally an $O_q(G)$-bimodule by 
\begin{equation}
\psi(\varphi\otimes a)\psi'
=\sum\psi\varphi\psi'_{(0)}\otimes a\psi'_{(1)}
\qquad(\varphi, \psi, \psi'\in O_q(G), a\in\wedge\m^*).
\end{equation}


By Lemma \ref{lem:adm} $\m$ is a left $\K$-module via the adjoint action.
Hence we have a right $\K$-module structure 
\begin{equation}
\m^*\otimes\K\to\m^*
\qquad(a\otimes k\mapsto a\al k)
\end{equation}
of $\m^*$ given by
\[
\langle a\al k,x\rangle=\langle a,\ad(k)(x)\rangle
\qquad(a,\in\m^*, k\in\K, x\in\m).
\]



\begin{lemma}
\label{lem:mK}
The right $\K$-module structure of $\m^*$ is naturally extended to that of $\wedge\m^*$ by
\begin{itemize}
\item[(a)]
$1\al k=\varepsilon(k)1$\quad$(k\in\K)$,
\item[(b)]
$(a\wedge b)\al k=\sum(a\al k_{(1)})\wedge(b\al k_{(0)})$
\quad
$(a, b\in\wedge\m^*, k\in\K)$.

\end{itemize}
\end{lemma}
We also use the left $\K$-module structure 
\begin{equation}
\K\otimes\wedge\m^*\to\wedge\m^*
\qquad(k\otimes a\mapsto k\ar a)
\end{equation}
of $\wedge\m^*$ given by
\[
k\ar a=a\al Sk
\qquad(k\in \K).
\]
We define a left $\K$-module structure 
\begin{equation}
\K\otimes(O_q(G)\otimes\wedge\m^*)\to O_q(G)\otimes\wedge\m^*
\qquad(k\otimes \omega\mapsto k\ar \omega)
\end{equation}
of $O_q(G)\otimes\wedge\m^*$ by
\[
k\ar(\varphi\otimes a)=\sum k_{(0)}\varphi\otimes (k_{(1)}\ar a).
\]
This gives a right $O_q(L_J)$-comodule structure of $O\otimes\wedge\m^*$.

The right $U_q(\Gg)$-module structure of 
$O_q(G)\otimes\wedge\m^*$
\begin{equation}
(\varphi\otimes a)u=
\varphi u\otimes a
\qquad(\varphi\in O_q(G), a\in\wedge\m^*, u\in U_q(\Gg))
\end{equation}
gives a left $O_q(G)$-comodule structure of 
$O_q(G)\otimes\wedge\m^*$.

For $n\geqq0$ we define a linear map
\begin{equation}
d:O_q(G)\otimes\wedge^n\m^*\to O_q(G)\otimes\wedge^{n+1}\m^*
\end{equation}
by 
\[
d(\varphi\otimes a)=\sum_{\beta\in\Xi_J}e_\beta\varphi\otimes (e^*_\beta\wedge a)
\qquad(\varphi\in O_q(G), a\in\wedge^n\m^*).
\]
\begin{theorem}
\label{thm:1}
\begin{itemize}
\item[(i)]
The linear endomorphism $d$ of $O_q(G)\otimes\wedge\m^*$ of degree 1 gives a dg-algbra structure of the graded algebra $O_q(G)\otimes\wedge\m^*$.
Namely, we have $d^2=0$ and 
\begin{align}
d(\omega\wedge \xi)=&d\omega\wedge\xi+(-1)^n \omega\wedge d\xi
\\
\nonumber
&\qquad(\omega\in O_q(G)\otimes\wedge^n\m^*,\; \xi\in
O_q(G)\otimes\wedge\m^*).
\end{align}
\item[(ii)]
The graded algebra $O_q(G)\otimes\wedge\m^*$ is generated by 
$O_q(G)=O_q(G)\otimes\wedge^0\m^*$ and 
$\Image(d:O_q(G)\otimes\wedge^0\m^*
\to O_q(G)\otimes\wedge^1\m^*)$.
Hence $O_q(G)\otimes\wedge\m^*$ is a differential calculus over $O_q(G)$ in the sense of \cite{BeM}.
\item[(iii)]
The differential calculus $O_q(G)\otimes\wedge\m^*$ is the maximal prolongation of the first order differential calculus 
$d:O_q(G)\to O_q(G)\otimes\m^*$.
Namely, for any differential calculus $\Omega$ which is a prolongation of the first order differential calculus 
$d:O_q(G)\to O_q(G)\otimes\m^*$ there exists uniquely a morphism of differential calculus $O_q(G)\otimes\wedge\Gm^*\to\Omega$ which is identity on the degree zero and one parts.

\item[(iv)] 
The differential calculus $O_q(G)\otimes\wedge\m^*$ is 
$(O_q(G),O_q(L))$-bicovariant.
Namely, the multuplication 
\[
(O_q(G)\otimes\wedge\m^*)\otimes (O_q(G)\otimes\wedge\m^*)\to O_q(G)\otimes\wedge\m^*
\]
and $d:O_q(G)\otimes\wedge\m^*\to O_q(G)\otimes\wedge\m^*$ are homomorphisms of $(O_q(G),O_q(L))$-bicomodules.
\end{itemize}
\end{theorem}
\subsection{}
We set
\[
\F=\{f\in O_q(G)\mid kf=\varepsilon(k)f\;(k\in\K)\}.
\]
We have
\[
\Delta(\F)\subset O_q(G)\otimes \F,
\]
so that $\F$ is a left coideal subalgebra of $O_q(G)$.

Since $O_q(G)\otimes\wedge\m^*$ is an $(O_q(G),O_q(L_J))$-bicovariant dg-algebra, we obtain a left $O_q(G)$-covariant dg-algebra
\[
\Omega=\bigoplus_{n\geqq0}\Omega^n,
\]
where
\begin{align*}
\Omega^n=&(O_q(G)\otimes\wedge^n\m^*)^{\co(O_q(L_J))}
\\
=&
\{\omega\in O_q(G)\otimes\wedge^n\m^*\mid
k\ar\omega=\varepsilon(k)\omega\;(k\in\K)\}.
\end{align*}
Regarding $\wedge\m^*$ as a left $O_q(L_J)$-comodule via
the right $\K$-module structure of $\wedge\m^*$ given in Lemma \ref{lem:mK} we can also write
\[
\Omega
=
O_q(G)\Box_{O_q(L_J)}\wedge\m^*.
\]
Here, for a coalgebra $C$, a right $C$-comodule $M$, a left $C$-comodule $N$ we denote by $M\Box_CN$ the cotensor product of $M$ and $N$.
We have 
\begin{align*}
\Omega^0=
O_q(G)\Box_{O_q(L_J)}\BC=\F,
\qquad
\Omega^1=
O_q(G)\Box_{O_q(L_J)}\m^*.
\end{align*}


\begin{theorem}
\label{thm:2}
\begin{itemize}
\item[(i)]
$\Omega$ is a left $O_q(G)$-covariant differential calculus over $\F$.
\item[(ii)]
$\Omega$ is the maximal prolongation of the first order differential calculus 
$d:\F\to O_q(G)\Box_{O_q(L_J)}\m^*$.
\item[(iii)]
The right $\F$-module structure of $\Omega/\F^+\Omega$ is trivial.
Namely, we have
$\omega f=\varepsilon(f)\omega$ for $\omega\in\Omega/\F^+\Omega$.
\end{itemize}
\end{theorem}
\section{Cohomology groups of vector bundles}
Let $V$ be an integrable right $U_q(\Gp_J)$-module.
We associate to it a left $O_q(G)$-covariant left $\F$-module
\[
\CV=O_q(G)\Box_{O_q(L_J)}V.
\]
The left $O_q(G)$-comodule structure and the left $\F$-module structure are given by
\[
(\sum\varphi_r\otimes v_r)u=
\sum\varphi_ru\otimes v_r,\qquad
f\sum\varphi_r\otimes v_r=\sum f\varphi_r\otimes v_r
\]
for $\sum\varphi_r\otimes v_r\in O_q(G)\Box_{O_q(L_J)}V$, 
$u\in U_q(\Gg)$, $f\in\F$.
Recall that $\Omega$ is a left $O_q(G)$-covariant $\F$-bimodule.
Hence we have a left $O_q(G)$-covariant left $\F$-module 
$\Omega\otimes_{\F}\CV$.
By Takeuchi's equvalence \cite{Tak} and Theorem \ref{thm:2} (iii) we see that
\[
\Omega\otimes_{\F}\CV\cong 
O_q(G)\Box_{O_q(L_J)}(\wedge\m^*\otimes V),
\]
where the left $O_q(L_J)$-comodule structure of 
$\wedge\m^*\otimes V$ is given by 
\[
(a\otimes v)k=\sum (a\al k_{(0)})\otimes v k_{(1)}
\qquad
(a\in \wedge\m^*, v\in V, k\in\K).
\]
We can rewrite it as 
\[
\Omega\otimes_{\F}\CV\cong 
(O_q(G)\otimes V\otimes(\wedge\m^*))^{\co(O_q(L_J))},
\]
where $O_q(G)\otimes V\otimes(\wedge\m^*)$ is an 
$(O_q(G),O_q(L_J))$-bicomodule 
given by
\[
(\varphi\otimes v\otimes a)u
=\varphi u\otimes v\otimes a,\quad
k(\varphi\otimes v\otimes a)
=\sum k_{(0)}\varphi\otimes v(Sk_{(1)})\otimes (k_{(2)}\ar a)
\]
for
$\varphi\in O_q(G)$, $v\in V$, $a\in\wedge\m^*$, $u\in U_q(\Gg)$, $k\in\K$.

Define
\[
\tilde{\nabla}_V:O_q(G)\otimes V\otimes \wedge\m^*
\to O_q(G)\otimes V\otimes \wedge\m^*
\]
by
\[
\tilde{\nabla}_V(\varphi\otimes v\otimes a)
=
\sum_{\beta\in\Xi_J}
e_{\beta(0)}\varphi\otimes v(Se_{\beta(1)})\otimes e^*_\beta\wedge a
\quad
(\varphi\in O_q(G), v\in V, a\in\wedge\m^*).
\]
Since $\tilde{\nabla}_V$ is an endomorphism of an $(O_q(G),O_q(L_J))$-bicomodule, it induces an endomorphism
\[
\nabla_V:\Omega\otimes_{\F}\CV\to\Omega\otimes_{\F}\CV
\]
of the left $O_q(G)$-comodule $\Omega\otimes_{\F}\CV$.
\begin{theorem}
%\label{thm:}
\begin{itemize}
\item[(i)]
We have
\[
\nabla_V(\omega\wedge s)=d\omega\wedge s+(-1)^n\omega\wedge\nabla_V(s)
\qquad(
\omega\in\Omega, s\in \Omega\otimes_{\F}\CV).
\]
\item[(ii)] We have $\nabla_V^2=0$.
\end{itemize}
\end{theorem}

For $r\geqq0$ we define a left $O_q(G)$-comodule
 $H^r(V)$ as the $r$-th cohomology group of the
complex
\[
\Omega^\bullet\otimes_{\F}\CV
=
[
0\to \Omega^0\otimes_{\F}\CV\xrightarrow{\nabla_V}\Omega^1\otimes_{\F}\CV\xrightarrow{\nabla_V}\
\cdots].
\]

Now we compare $H^r(V)$ with the derived induction functor of \cite{APW}.
Consider the  left exact functor 
\[
\Ind:{}^{O_q(P_J)}\!\Mod\to {}^{O_q(G)}\!\Mod
\]
given by
\[
\Ind(V)=O_q(G)\Box_{O_q(P_J)}V.
\]
Then we have
\begin{align*}
\Ind(V)\cong&(O_q(G)\otimes V)^{\co(O_q(P_J))}
\\
\cong&
\Hom_{U_q(\Gp_J)}(\BC,O_q(G)\otimes V)
\\
\cong&
\Hom_{U_q(\Gm_J)}(\BC,O_q(G)\otimes V)^{\co(O_q(L_J))},
\end{align*}
where $O_q(G)\otimes V$ is regarded as a right $O_q(P_J)$-comodule by
\[
y(\varphi\otimes v)=\sum y_{(0)}\varphi\otimes v(Sy_{(1)})
\qquad (\varphi\in O_q(G), v\in V, y\in U_q(\Gp_J)).
\]
By our assumption on $q$ the category $\Mod^{O_q(L_J)}$ is semisimple, and hence the derived functor of $\Ind$ is given by
\[
R^r\Ind(V)\cong\Ext^r_{U_q(\Gm_J)}(\BC,O_q(G)\otimes V)^{\co(O_q(L_J))}.
\]
By a general result on quadratic algebras (see \cite{P}, \cite{PP}) we have a free resolution 
\[
\cdots
\xrightarrow{\deru}
U_q(\Gm_J)\otimes(\wedge^1\m^*)^*
\xrightarrow{\deru}
U_q(\Gm_J)\otimes(\wedge^0\m^*)^*
\to
\BC\to 0
\]
of the trivial $U_q(\Gm_J)$-module $\BC$,
where
\[
\deru:U_q(\Gm_J)\otimes(\wedge^n\m^*)^*
\to
U_q(\Gm_J)\otimes(\wedge^{n-1}\m^*)^*
\]
is given by 
\[
\deru(u\otimes x)=\sum_{\beta\in\Xi_J}ue_\beta\otimes xe^*_\beta
\qquad(u\in U_q(\Gm_J), x\in(\wedge^n\m^*)^*).
\]
Here, the right $\wedge\m^*$-module structure of $(\wedge\m^*)^*$ is given by
\[
\langle xa,b\rangle=\langle x, a\wedge b\rangle
\qquad(x\in (\wedge\m^*)^*, a, b\in \wedge\m^*).
\]
Hence $R^r\Ind(V)$ is given by the $r$-th cohomology group of the complex
\begin{align*}
&\Hom_{U_q(\Gm_J)}(U_q(\Gm_J)\otimes(\wedge^\bullet\m^*)^*,O_q(G)\otimes V)^{\co(O_q(L_J))}
\\
\cong&
\Hom_{\BC}((\wedge^\bullet\m^*)^*,O_q(G)\otimes V)^{\co(O_q(L_J))}
\\
\cong&
(O_q(G)\otimes V\otimes(\wedge^\bullet\m^*))^{\co(O_q(L_J))}.
\end{align*}
It is easily checked that this complex is isomorphic to $\Omega^\bullet\otimes_{\F}\CV$.
Therefore, we obtain the following.
\begin{theorem}
%\label{thm:}
For a $U_q(\Gp_J)$-module $V$ we have $H^r(V)\cong R^r\Ind(V)$.
\end{theorem}
%\section*{Acknowledgment}
\bibliographystyle{unsrt}
\begin{thebibliography}{99}
%
\bibitem{APW}
Andersen, H. H.;
Polo, P.;
Wen, K.
Representations of quantum algebras. 
Invent. Math. {\textbf{104}} (1991), no. 1, 1--59.

\bibitem{BeM}
Beggs, Edwin J.; Majid, Shahn.
Quantum Riemannian geometry,
Grundlehren Math. Wiss., 355[Fundamental Principles of Mathematical Sciences]
Springer, Cham, 2020, xvi+809 pp.

\bibitem{BeS}
Beggs, Edwin; Paul Smith, S. Non-commutative complex differential geometry. J. Geom. Phys. 72 (2013), 7--33. 

\bibitem{CDO}
Carotenuto, Alessandro; D\'{i}az Garc\'{i}a, Fredy; \'{O} Buachalla, Reamonn,
A Borel-Weil theorem for the irreducible quantum flag manifolds
Int. Math. Res. Not. IMRN 2023, no. 15, 12977--13006. 

\bibitem{HK1}
Heckenberger, Istv\'{a}n; Kolb, Stefan. 
Differential calculus on quantum homogeneous spaces. Lett. Math. Phys. 63 (2003), no. 3, 255--264.

\bibitem{HK2}
Heckenberger, Istv\'{a}n; Kolb, Stefan. 
The locally finite part of the dual coalgebra of quantized irreducible flag manifolds. 
Proc. London Math. Soc. (3) 89 (2004), no. 2, 457--484. 

\bibitem{HK3}
Heckenberger, Istv\'{a}n; Kolb, Stefan. De Rham complex for quantized irreducible flag manifolds. J. Algebra 305 (2006), no. 2, 704--741.

\bibitem{KMT}
Kamita, Atsushi; Morita, Yoshiyuki; Tanisaki, Toshiyuki. 
Quantum deformations of certain prehomogeneous vector spaces. I. Hiroshima Math. J. 28 (1998), no. 3, 527--540.  

\bibitem{Lb}
Lusztig, G.
Introduction to quantum groups. 
Progr. Math., \textbf{110}, Boston etc. Birk\"hauser, 1993.



\bibitem{PP}
Polishchuk, Alexander; Positselski, Leonid. 
Quadratic algebras. University Lecture Series, 37. American Mathematical Society, Providence, RI, 2005. xii+159 pp.

\bibitem{P}
Priddy, Stewart B. 
Koszul resolutions. Trans. Amer. Math. Soc. 152 (1970), 39--60.



\bibitem{Se}
Serre, Jean-Pierre.
G\'{e}ométrie alg\'{e}brique et g\'{e}ométrie analytique.
Ann. Inst. Fourier (Grenoble) 6 (1955/56), 1--42.
\bibitem{Tak}
Takeuchi, Mitsuhiro. 
Relative Hopf modules--equivalences and freeness criteria. 
J. Algebra 60 (1979), no. 2, 452--471.  

\bibitem{TM}
Tanisaki, Toshiyuki. Modules over quantized coordinate algebras and PBW-bases. J. Math. Soc. Japan 69 (2017), no. 3, 1105--1156. 



\end{thebibliography}
\end{document}